% Nouveau EGA V (cote 25): archive cote 25, site page 130 (PDF page 131).
% English translation of the checked French annotated master.
% AI-translated by GPT-6 Astra (OpenAI, ChatGPT Pro mode).
\scanpage{25/130}
\mspage{\hand{33}}

which reduces us to the case where S is local artinian, S=Spec(A). One proceeds
by induction on the integer n such that $m^{n+1}=0$ , which allows us to
assume that the two sections under consideration \var{\inserted[typed]{u,v}}{u,v}{insertion}{ega25_25_130-v0001} are equal mod $m^{n\var{\xout{\unsure{+1}}}{}{typewriter strike-out}{ega25_25_130-v0002}}$. Thus
one is obtained from the other by means of an element \var{\xout{xxxxxx(xxxx xxxxxxx}}{}{typewriter strike-out}{ega25_25_130-v0003} \hand{$\delta$ of}
$\mathrm{Hom}_k(u_o^{\hand{*}}(\hand{\Omega}^1_{F_o/k}), V)$ , where $k=A/m$ is the residue field, $F_o=F\otimes_A k$ is the
reduced fiber, $V=m^n = m^n/m^{n+1}$ . One must prove that \hand{$\delta$}$= 0$, under
the hypothesis $h(u)=h(v)$ . The general principle of the verification is as
follows: first expressing that $h(u)$ and $h(v)$ coincide mod $m^n$, one sees
that their ``difference'' can be written as an element \hand{$\delta$}$'$ of
$\mathrm{Hom}_k(w_o^{\hand{*}}(\hand{\Omega}^1_{G_o/k}),V)$ , where $w_o = h_o(u_o)=h_o(v_o)$ \var{\xout{xxx}}{}{typewriter strike-out}{ega25_25_130-v0004} and where $G_o = G\otimes_A k$ ;
this element is precisely the on\uline{e} obtained from \var{\inserted[typed]{\hand{$\delta$} by composition with}}{\hand{$\delta$} by composition with}{insertion}{ega25_25_130-v0005} the natural
homomorphism $h_o^{\hand{*}} : w_o^{\hand{*}}(\hand{\Omega}^1_{G_o/k}) \hand{\to} u_o^{\hand{*}}(\hand{\Omega}^1_{F_o/k})$ obtained from $h_o: F_o\hand{\to} G_o$. Since
$h(u)=h(v)$, hence \hand{$\delta$}$' = 0$, the composite of \hand{$\delta$} with the preceding homomorphism
$h_o^{\hand{*}}$ is zero, so knowing that $h_o$ is surjective implies that \hand{$\delta$}$=0$
and we win. Now the fact that $h_o:F_o\hand{\to} G_o$ is a monomorphism, hence induces
a monomorphism on the set of points with values in the dual
numbers over k, does indeed imply that $h_o^{\hand{*}}$ is surjective (its transpose
being injective). This argument is valid as it stands when G
is representable, which is however not the case in the case that
interests us. \var{\deleted{One must therefore}}{}{deletion}{ega25_25_130-v0006} \var{\inserted{One can nevertheless still}}{One can nevertheless still}{insertion}{ega25_25_130-v0007} define a vector space \hand{$\mathcal{V}_{w_o}$} \var{\inserted{(tangent space to $G_o/k$ at $w_o$)}}{(tangent space to $G_o/k$ at $w_o$)}{insertion}{ega25_25_130-v0008} over k, playing the role
of the \var{\inserted{dual of}}{dual of}{insertion}{ega25_25_130-v0009} $w_o^{\hand{*}}(\hand{\Omega}^1_{G_o/k})$ in expressing the ``deviation'' of two points of G that
coincide mod $m^n$ \var{\inserted{as an element of $\mathcal{V}_{w_o}\otimes_k V$}}{as an element of $\mathcal{V}_{w_o}\otimes_k V$}{insertion}{ega25_25_130-v0010}. This is essentially straightforward, and contained
in the systematic developments of par. 2\var{\over{5}{6}}{6}{writing over}{ega25_25_130-v0011} (? Infinitesimal
extensions), to which one could refer here. In case a), G being the
functor $\mathrm{Div}_{X/S}$ , $w_o$ thus corresponds to a Cartier divisor \var{\xout{xx}}{}{typewriter strike-out}{ega25_25_130-v0012} $D_o$ on
$X_o = X\otimes_A k$ , and one must take \hand{$\Omega$}$= H^o(D_o, \underline{n}_{D_o/X_o})$ where $\underline{n}$ is the
normal sheaf of $D_o$ in $X_o$ , also isomorphic to the sheaf induced
on $D_o$

\mspage{\hand{130}}
\tnote{25-130-1}{The final-letter underline in French ``celui'' is carried by the final letter of its English rendering ``the one''. The word ``straightforward'' is already English in the manuscript. The source says that the displayed morphism h with subscript o is surjective at the indicated step; its notation is preserved without inserting a star.}
